\section {Introduction}

\subsection{Overview of Quadtrees}
A \textbf{Quadtree} is a tree data structure in which each internal node has exactly four children. It is most often used to partition a two-dimensional space by recursively subdividing it into four quadrants or regions. The data represented may be points, regions, lines, or curves. Quadtrees are fundamental in fields that require spatial indexing, such as computer graphics, geographic information systems (GIS), and image processing. By using a hierarchical representation, Quadtrees allow for efficient spatial queries, reducing the time complexity of operations like range searching or collision detection from $O(N)$ to $O(\log N)$.

\subsection{Historical Background}
The term "Quadtree" was first coined by \textbf{Raphael Finkel} and \textbf{J.L. Bentley} in \textbf{1974} to describe a data structure for retrieving points based on their coordinates in a 2D plane. However, the concept of recursive decomposition of space has older roots in mathematics and cartography. 

Shortly after its introduction, the structure was adapted into the \textbf{Region Quadtree} by \textbf{Hunter} in \textbf{1978} for image processing and computer graphics. During the 1980s and 1990s, new variants such as the \textit{Point Quadtree}, \textit{MX Quadtree}, and \textit{PR Quadtree} were developed for different data types. Today, Quadtrees power many practical systems, including map tiling and collision detection.


\clearpage